\documentclass[12pt, a4paper]{article}

\usepackage[utf8]{inputenc}
\usepackage[T1]{fontenc}
\usepackage{geometry}
\usepackage{hyperref}
\usepackage{graphicx}
\usepackage{natbib}
\usepackage{amsmath}
\usepackage{amssymb}
\usepackage{booktabs}
\usepackage{xcolor}

\geometry{left=25mm, right=25mm, top=25mm, bottom=25mm}

\title{\textbf{Where the Signal Is:\\
Label Provenance, Encoder Choice and the Limits of Fusion\\
in Arabic Multimodal Affect Recognition}}
\author{Nour Gaser, Mohamed Wael \\
\textit{Department of Computer Science} \\
\textit{German International University}}
\date{August 2026}

\begin{document}
\maketitle

\begin{abstract}
Cross-modal attention is the dominant architectural response to multimodal
affect recognition. We set out to test whether it outperforms late fusion for
Arabic, and in particular whether its advantage concentrates on samples where
the modalities carry conflicting evidence. Across two corpora, twelve fusion
arrangements and four encoder families, it does not --- and the reason is
instructive. On Ar-MuSA, the only openly available Arabic multimodal sentiment
corpus, we find that the overall label agrees with the corpus's own
transcript-level annotation on \textbf{96.5\%} of samples, so a text-only model
already holds nearly all the label information and no fusion mechanism has room
to distinguish itself. We predicted from this that a corpus whose label is
\emph{not} recoverable from text would behave differently, and tested that
prediction on EGY-MER, an Egyptian Arabic emotion corpus obtained during this
work. The prediction holds: on EGY-MER the same measurement gives
\textbf{59.4\%} for transcript against \textbf{91.1\%} for audio, the ordering
of the modalities inverts, and the multimodal gain over the best single modality
triples. Yet even there, no attention arrangement we tried beats plain
concatenation by more than the variation between seeds, and concatenation is the
smallest model by a factor of two to seven. Every measurable
improvement we obtained came from representation --- dialect-matched text
encoders ($+0.054$), reading audio from mid-stack rather than the final layer
($+0.053$), expression-pretrained visual encoders ($+0.026$) and fine-tuning
($+0.051$) --- and none from fusion design. We report four falsified hypotheses
alongside these. Our system reaches weighted $F_1$ 0.664 on EGY-MER against a
published baseline of 0.519 (single seed; the three-seed mean is 0.654, and 0.669
with the best encoders).
\end{abstract}

\section{Introduction}
\label{sec:intro}

Multimodal affect recognition rests on the premise that non-verbal channels
carry information the words do not. The canonical motivating case is
incongruence: an utterance whose literal content is contradicted by its
delivery. Cross-modal attention \citep{Tsai2019} operationalises this directly,
letting one modality condition on another before a decision is made, and has
become the default architectural response; it has recently reached Arabic
\citep{Cheikhi2026}.

This report asks a narrow architectural question and arrives at a broader answer
about evaluation. Section~\ref{sec:setup} describes the two corpora and the
protocol. Section~\ref{sec:provenance} introduces \emph{label provenance}, the
measurement on which everything else turns, and reports it for both corpora.
Section~\ref{sec:fusion} reports the fusion comparison.
Section~\ref{sec:encoders} reports the encoder study, which is where the actual
gains are. Section~\ref{sec:falsified} lists the hypotheses we tested and
rejected. Section~\ref{sec:limits} states the limitations.

\section{Corpora and Protocol}
\label{sec:setup}

\paragraph{Ar-MuSA} \citep{Khaled2025} comprises 8{,}700 Egyptian Arabic and MSA
clips (8{,}699 usable) labelled for three-way sentiment with a binary sarcasm
flag. Crucially, both are annotated \emph{separately for the transcript, the
audio and the frame}, in addition to an overall judgement. The visual stream is
a single extracted frame per clip.

\paragraph{EGY-MER} \citep{Hussein2025} is an Egyptian Arabic emotion corpus,
obtained from the corpus holders during this work. As delivered it holds 5{,}471
speaker-diarised segments with media, of which \textbf{505 carry a complete
annotation} --- three annotators, three modalities each, with a documented
majority-vote resolution --- across 49 source videos and 70 minutes. We use the
484 whose overall label is one of the seven emotions. Unlike Ar-MuSA it ships
full video segments, so the visual stream is genuinely temporal; we sample 16
frames per segment.

\paragraph{Protocol.} Encoders are frozen and features extracted once, so that
the only difference between compared models is the fusion mechanism.
Validation is stratified $k$-fold with identical folds for every model, and
predictions are collected out-of-fold so subsets can be analysed without
leakage. Because EGY-MER has 484 samples across 7 classes, every EGY-MER number
is a mean over three seeds with the standard deviation reported; we treat
differences smaller than the seed spread as ties rather than results.

\section{Label Provenance}
\label{sec:provenance}

Both corpora annotate each modality separately, so we can ask how much of the
overall label a single modality's own annotation already supplies. We call this
\emph{label provenance}, and it turns out to determine whether a corpus can
evaluate fusion at all.

\begin{table}[h]
\centering
\small
\begin{tabular}{lcc}
\toprule
\textbf{Annotation} & \textbf{Ar-MuSA (sentiment)} & \textbf{EGY-MER (emotion)} \\
\midrule
Transcript & \textbf{96.5\%} & 59.4\% \\
Audio      & 77.1\% & \textbf{91.1\%} \\
Frame / video & 48.2\% & 84.4\% \\
\midrule
Overall label differs from text label & 3.5\% & \textbf{40.6\%} \\
\bottomrule
\end{tabular}
\caption{How often the overall label equals each modality's own annotation.
The ordering inverts between the two corpora. Computed over samples carrying all
three modality labels (8{,}699 and 473) by \texttt{experiments/label\_provenance.py}.}
\label{tab:provenance}
\end{table}

On Ar-MuSA a predictor with perfect access to the transcript label alone would
score macro-$F_1$ 0.9655 on the ostensibly multimodal target; only 304 of 8{,}699
clips disagree with the text. The sarcasm flag behaves the same way (94.9\%).
Any fusion comparison on that corpus is therefore a contest over a
$\sim$3.5\% residue, and a model with no access to audio or vision begins with
almost all the available label information. This is a property of the benchmark
and applies to every result reported on it, ours included.

EGY-MER inverts this. Its label lives in the \emph{delivery}: text is the
weakest channel, audio the strongest, and 40.6\% of segments carry an overall
label that differs from the text label. This is precisely the regime in which
non-verbal channels must carry discriminative information, and it is what the
Ar-MuSA result predicted would be required.

\paragraph{The prediction is confirmed by model behaviour.} Unimodal macro-$F_1$
inverts to match the annotations (Table~\ref{tab:unimodal}), and the gain from
using all three modalities over the best single one is \textbf{three times
larger} on EGY-MER ($+0.0439$) than on Ar-MuSA ($+0.0135$).

\begin{table}[h]
\centering
\begin{tabular}{lcc}
\toprule
\textbf{Modality alone} & \textbf{Ar-MuSA} & \textbf{EGY-MER} \\
\midrule
Text   & \textbf{0.5458} & 0.3404 \\
Audio  & 0.4184 & 0.4770 \\
Visual & 0.4327 & \textbf{0.5225} \\
\midrule
All three (late fusion) & 0.5593 & 0.5664 \\
Multimodal gain & $+0.0135$ & $\mathbf{+0.0439}$ \\
\bottomrule
\end{tabular}
\caption{Unimodal macro-$F_1$. The modality ordering inverts with the task, in
the same direction as label provenance.}
\label{tab:unimodal}
\end{table}

\section{Fusion Arrangements}
\label{sec:fusion}

On EGY-MER --- the corpus where multimodality demonstrably matters --- we
compared twelve fusion arrangements under one protocol: plain concatenation; a
capacity-matched widening of it; per-modality gating (the strategy the published
EGY-MER baseline uses); cross-modal attention anchored on each of the three
modalities in turn; an all-pairs arrangement in which every modality attends to
every other (following \citealp{Tsai2019}); full self-attention over the
concatenated token sequence; and a hybrid described below.

\textbf{With the base encoders, plain concatenation scores highest and is the
smallest model} (891K parameters against 1.1--6.4M). The largest seed standard
deviation is 0.021, so the top seven arrangements are statistically
indistinguishable; the remaining five are worse. Repeating the benchmark with the
best encoders from Section~\ref{sec:encoders} changes the order at the top: a
gated hybrid scores 0.608 against concatenation's 0.594, a difference smaller than
its own seed deviation (0.017). In neither run does any arrangement beat
concatenation by more than seed variation, and the three Tri-CAN variants fall in
the lower half both times.

\paragraph{A design flaw in the proposed architecture.} The originally specified
Tri-CAN feeds its classifier
$[\,\text{pool}(\text{anchor}),\ \text{pool}(\text{anchor's view of }A),\
\text{pool}(\text{anchor's view of }B)\,]$, where every term is pooled over
\emph{anchor} positions. All three therefore live in the anchor's token space,
and the raw $A$ and $B$ representations never reach the classifier. Tri-CAN is
not late fusion plus attention: it \emph{replaces} the other modalities with the
anchor's interrogation of them, and so has strictly less information than the
baseline it is meant to beat. This predicts the observed pattern. On
incongruent samples every anchor scores below late fusion ($-0.047$, $-0.070$ and
$-0.066$ for text, audio and visual anchors), while on congruent samples the
text-anchored model is ahead ($+0.036$). The difference-in-differences for the
text anchor is $-0.082$, but its 95\% bootstrap interval ($[-0.189, +0.034]$,
208 and 265 samples) includes zero, so we read it as a consistent direction
rather than an established effect. On a corpus that is 56\% incongruent, the
losses dominate.

Restoring the discarded representations (our \emph{hybrid} arrangement, whose
input is a strict superset of late fusion's, here with a learned gate) lifts
text-anchored attention from
$-0.0099$ to $-0.0014$ relative to the baseline --- from losing to exactly tied.
The diagnosis is therefore confirmed, but the repair only recovers parity: the
attended view adds nothing the raw vectors did not already carry.

\section{Where the Gains Actually Are}
\label{sec:encoders}

Every measurable improvement in this study came from representation, not
architecture:

\begin{table}[h]
\centering
\small
\begin{tabular}{p{0.56\linewidth}cc}
\toprule
\textbf{Change} & \textbf{Gain (macro-$F_1$)} & \textbf{Corpus} \\
\midrule
Dialect-matched text encoder (AraBERT-twitter over AraBERTv2) & $+0.054$ & Ar-MuSA \\
Fine-tuning the text encoder & $+0.051$ & Ar-MuSA \\
Reading audio from mid-stack rather than the final layer & $+0.053$ & EGY-MER \\
Expression-pretrained visual encoder over ImageNet Swin & $+0.026$ & EGY-MER \\
\midrule
Best fusion architecture over plain concatenation & $\approx 0$ & both \\
\bottomrule
\end{tabular}
\caption{Representation choices dominate architectural ones by an order of
magnitude.}
\label{tab:gains}
\end{table}

The audio result deserves emphasis because it invalidates a conclusion we had
already drawn. In wav2vec2-style models the final layers specialise toward the
self-supervised pretraining objective, discarding the slowly-varying
paralinguistic content an affect task needs. Reading the final layer cost
$0.045$--$0.053$ macro-$F_1$, and our earlier finding that audio ``contributes
almost nothing'' was measured at the wrong depth.

Fine-tuning is separable from fusion: on Ar-MuSA it improved late fusion and
Tri-CAN by exactly $+0.0505$ each, leaving their ordering unchanged. The fusion
null is therefore not an artefact of the frozen-feature protocol.

\section{Hypotheses We Rejected}
\label{sec:falsified}

We state these because each was predicted in advance and each failed.

\begin{enumerate}
\item \textbf{Anchor choice.} Since text is EGY-MER's weakest modality, we
predicted anchoring attention on audio or vision would help. It did not: text
remained the best anchor (audio $-0.034$, vision $-0.009$ relative to it).
\item \textbf{Multilingual audio encoder.} XLS-R-300m scored \emph{worse} than
English-pretrained wav2vec2-base ($0.3835$ vs $0.4206$), so the weak audio
signal was not a language mismatch.
\item \textbf{Face cropping.} Cropping to the detected face before an
expression-pretrained encoder cost $-0.0375$ despite 98\% detection. We
attribute this to the corpus being speaker-diarised but not
speaker-\emph{localised}: the largest face in a talk-show frame is not
necessarily the person speaking, so cropping reliably delivers the wrong face.
Active-speaker detection would be required.
\item \textbf{Emotion-pretrained audio encoder.} A model fine-tuned for speech
emotion recognition scored worse than the generic self-supervised one
($0.4578$ vs $0.4770$); it is tuned on English acted emotion, and that
specialisation narrows rather than improves the representation for Egyptian
Arabic.
\end{enumerate}

Items 2 and 4 share a lesson worth stating: matching an encoder's \emph{domain}
to the data helps, while matching its \emph{task label space} does not.

\section{Limitations}
\label{sec:limits}

\begin{enumerate}
\item \textbf{Sample size.} EGY-MER's 484 annotated segments across seven
classes, with Fear at 12 and Happiness at 18, make macro-$F_1$ noisy. We report
mean $\pm$ SD over seeds and decline to interpret differences below the seed
spread.
\item \textbf{Baseline comparability.} Our weighted $F_1$ 0.664 (single seed;
0.654 over three) exceeds the
published 0.519, but that figure was obtained on 363 samples with different
splits; the comparison is indicative, not like-for-like.
\item \textbf{Frozen non-verbal encoders.} Only the text encoder was fine-tuned,
a compute limitation rather than a design choice.
\item \textbf{Annotations, not signal.} ``The modalities disagree'' means the
annotators assigned different labels per stream. It is a good proxy for
conflicting evidence, but a proxy.
\item \textbf{Grouping.} Folds are stratified by label but not grouped by source
video, so some speaker leakage across folds is possible. This affects absolute
scores more than the within-study contrasts, which apply equally to every model.
\end{enumerate}

\section{Conclusion}

Cross-modal attention did not outperform late fusion on either Arabic corpus,
under any anchor, arrangement, encoder family, or training protocol we tested.
We do not read this as evidence that the mechanism is worthless. We read it as
two separate findings.

First, \textbf{a benchmark can be structurally incapable of evaluating fusion},
and the standard Arabic multimodal sentiment corpus is: with 96.5\% of its label
recoverable from the transcript annotation alone, there is no headroom in which
mechanisms could be distinguished. Label provenance is cheap to compute, and we
suggest reporting it whenever a multimodal benchmark is proposed.

Second, on a corpus that \emph{does} have the headroom, \textbf{the binding
constraint is representation, not fusion}. Encoder domain, layer depth and
fine-tuning moved results by $0.03$--$0.05$ macro-$F_1$ each; the best of twelve
fusion architectures moved them by approximately zero, and the simplest
arrangement was also the smallest and the best.

\bibliographystyle{plainnat}
\bibliography{../../rp1/report/references}

\end{document}
